\documentclass[11pt]{article}

\usepackage[margin=1in]{geometry}
\usepackage{amsmath,amssymb,amsthm,mathtools,microtype}
\usepackage[T1]{fontenc}
\usepackage[hidelinks]{hyperref}

\newtheorem{theorem}{Theorem}
\newtheorem{lemma}[theorem]{Lemma}
\newtheorem{corollary}[theorem]{Corollary}

\newcommand{\C}{\mathbb C}
\newcommand{\Eop}{\mathcal E}
\newcommand{\M}{\mathcal M}

\title{An Eight-Term Counterexample to the Special Image Conjecture
  in Three Pairs}
\author{Roy van Rijn\\{\small Independent researcher}}
\date{28 July 2026}
\hypersetup{
  pdftitle={An Eight-Term Counterexample to the Special Image Conjecture in Three Pairs},
  pdfauthor={Roy van Rijn},
  pdfsubject={An explicit three-pair Image-Mathieu counterexample},
  pdfkeywords={Image Conjecture, Special Image Conjecture,
    Mathieu subspace, Gaussian moments, contraction}
}

\begin{document}
\maketitle

\begin{abstract}
We give an explicit counterexample to the Special Image Conjecture in
three contraction pairs.  In
\(\C[\tau,w,v,t,z,y]\), let
\[
 f=\tau^3(t+z)\bigl(wt^3-vy(t+y)^2\bigr),\qquad g=z.
\]
For the standard contraction map \(\Eop\), we prove for every \(m\geq1\)
that
\[
 \Eop(f^m)=0,\qquad [t]\Eop(gf^m)=(3m+1)!\,m!.
\]
Thus every positive power of \(f\) lies in the image of the operators
\(\partial_t-\tau,\partial_z-w,\partial_y-v\), whereas no \(gf^m\)
does.  The polynomial \(f\) has eight terms and the multiplier is linear.
Consequently the minimum number of pairs in which the Special Image
Conjecture can fail is either two or three.
\end{abstract}

\section{The counterexample}

For \(r\geq1\), write
\[
 R_r=\C[\zeta_1,\ldots,\zeta_r,z_1,\ldots,z_r],
 \qquad
 \M_r=\sum_{i=1}^r(\partial_{z_i}-\zeta_i)R_r.
\]
The Special Image Conjecture asserts that \(\M_r\) is a Mathieu
subspace: if \(h^m\in\M_r\) for every \(m\geq1\), then
\(qh^m\in\M_r\) for every \(q\in R_r\) and all sufficiently large
\(m\); see Zhao~\cite{Zhao2010} and van den Essen, Wright, and
Zhao~\cite{EWZ2011}.

Define the contraction map
\[
 \Eop_r(\zeta^\alpha z^\beta)=\partial_z^\alpha z^\beta.
\]

\begin{lemma}\label{lem:kernel}
For every \(r\geq1\), one has \(\M_r=\ker\Eop_r\).
\end{lemma}

\begin{proof}
For \(h\in R_r\),
\(\Eop_r((\partial_{z_i}-\zeta_i)h)=0\), so
\(\M_r\subseteq\ker\Eop_r\).  Conversely,
\[
 \zeta_i h\equiv\partial_{z_i}h\pmod{\M_r}.
\]
Repeatedly applying this congruence to each monomial gives
\(p\equiv\Eop_r(p)\pmod{\M_r}\).  Since \(\Eop_r\) is the identity on
\(\C[z_1,\ldots,z_r]\), \(\Eop_r(p)=0\) implies \(p\in\M_r\).
\end{proof}

Use the pairs
\[
 (\tau,t),\qquad(w,z),\qquad(v,y)
\]
and put
\[
 \M_3=(\partial_t-\tau)R_3+
       (\partial_z-w)R_3+
       (\partial_y-v)R_3.
\]

\begin{theorem}\label{thm:main}
Let
\[
 \boxed{\quad
 f=\tau^3(t+z)\bigl(wt^3-vy(t+y)^2\bigr),\qquad g=z.
 \quad}
\]
Then, for every \(m\geq1\),
\[
 \Eop_3(f^m)=0,\qquad
 [t]\Eop_3(gf^m)=(3m+1)!\,m!.
\]
Consequently \(f^m\in\M_3\) and \(gf^m\notin\M_3\) for every
\(m\geq1\).
\end{theorem}

\section{Proof}

First omit the pair \((\tau,t)\).  On
\(\C[W,Z,V,Y]\), define
\[
 \mathcal F(W^aZ^bV^cY^d)
 =\delta_{a,b}\delta_{c,d}\,a!\,c!
\]
and set
\[
 P=(1+Z)\bigl(W-VY(1+Y)^2\bigr).
\]
Expanding according to the number \(k\) of factors chosen from the
second summand gives
\[
 P^m=\sum_{k=0}^m(-1)^k\binom mk
 W^{m-k}V^k(1+Z)^mY^k(1+Y)^{2k}.
\]
For a nonzero contraction, the \(W,Z\) pair selects the coefficient
of \(Z^{m-k}\), while the \(V,Y\) pair selects the constant term of
\((1+Y)^{2k}\).  Hence
\begin{equation}\label{eq:pure}
 \mathcal F(P^m)
 =m!\sum_{k=0}^m(-1)^k\binom mk=0.
\end{equation}
After multiplication by \(Z\), the first pair instead selects
\([Z^{m-k-1}](1+Z)^m=\binom m{k+1}\).  Therefore
\begin{equation}\label{eq:mixed}
 \mathcal F(ZP^m)
 =m!\sum_{k=0}^{m-1}(-1)^k\binom m{k+1}=m!.
\end{equation}

Homogenize \(P\) to total coordinate degree four:
\[
 \widetilde P
 =(t+z)\bigl(wt^3-vy(t+y)^2\bigr).
\]
Every term of \(P^m\) that survives \(\mathcal F\) has total
\((Z,Y)\)-degree \(m\).  Its homogenization therefore contains
\(t^{4m-m}=t^{3m}\), exactly matching the power \(\tau^{3m}\).
The new contraction contributes \((3m)!\), and
\eqref{eq:pure} gives
\[
 \Eop_3(f^m)=(3m)!\mathcal F(P^m)=0.
\]

For \(ZP^m\), a surviving term of \(P^m\) has total
\((Z,Y)\)-degree \(m-1\).  Its homogenization contains
\(t^{4m-(m-1)}=t^{3m+1}\).  Contracting against \(\tau^{3m}\)
leaves one power of \(t\) and contributes \((3m+1)!\).  By
\eqref{eq:mixed},
\[
 [t]\Eop_3(zf^m)
 =(3m+1)!\mathcal F(ZP^m)
 =(3m+1)!\,m!\neq0.
\]
Lemma~\ref{lem:kernel} completes the proof of
Theorem~\ref{thm:main}.

\begin{corollary}
If \(r_{\mathrm{SIC}}\) denotes the minimum number of contraction pairs
in which the Special Image Conjecture fails, then
\[
 2\leq r_{\mathrm{SIC}}\leq3.
\]
\end{corollary}

\begin{proof}
The upper bound is Theorem~\ref{thm:main}.  The one-pair case follows
from~\cite[Theorem 2.8]{EWZ2011}, applied to
\(A=\C[\tau]\) and the radical principal ideal \((\tau)\).
\end{proof}

\section{Size, scope, and verification}

The expanded polynomial \(f\) has eight terms, coefficients in
\(\{1,-1,-2\}\), bidegree \((4,4)\), and ordinary degree eight.
The multiplier \(g=z\) has bidegree \((0,1)\).

A five-pair counterexample communicated by Dvorsky~\cite{Dvorsky2026},
from Long's \(SU(2)\) construction~\cite{Long2026}, is
\[
 \zeta_t(\zeta_a\zeta_d-\zeta_b\zeta_c)
 (t+c)(ad+bt),\qquad g=-c.
\]
Theorem~\ref{thm:main} lowers this presently recorded five-pair
construction to three pairs while retaining eight terms and a linear
coordinate multiplier.  The new \(f\) is not of the separated form
\(\lambda(\zeta)P(z)\); consequently the example does not give a
three-variable Generalized Vanishing Conjecture counterexample.  No
claim concerning the two-pair case or absolute literature priority is
made.

The all-order proof above is independent of computation.  An exact
sparse-arithmetic checker expands \(f\), verifies the contractions
through \(m=10\), and checks the two binomial identities through
\(m=99\).  It is available at
\begin{center}
\small
\url{https://github.com/royvanrijn/jacobian-research/blob/main/scripts/verify_three_pair_image_mathieu_counterexample.py}.
\end{center}
Its SHA-256 digest for this version is
\begin{center}
\ttfamily\small
ac5408c1278e6fe2d3f78b408bd2165454833db42a669cd929570ba1f2e30d5b.
\end{center}
As a low-order checksum,
\[
 \Eop_3(f)=0,\qquad \Eop_3(zf)=24t+6z.
\]

\paragraph{Acknowledgment.}
Computational and editorial assistance from OpenAI Codex was used in
searching for, checking, and presenting the example.  The author takes
responsibility for the mathematical statements and proof.

\begin{thebibliography}{9}
\small

\bibitem{Zhao2010}
W.~Zhao,
\emph{Images of commuting differential operators of order one with
constant leading coefficients},
J. Algebra \textbf{324} (2010), 231--247.
\href{https://doi.org/10.1016/j.jalgebra.2010.04.022}
{doi:10.1016/j.jalgebra.2010.04.022}.

\bibitem{EWZ2011}
A.~van den Essen, D.~Wright, and W.~Zhao,
\emph{On the Image Conjecture},
J. Algebra \textbf{340} (2011), 211--224.
\href{https://doi.org/10.1016/j.jalgebra.2011.04.036}
{doi:10.1016/j.jalgebra.2011.04.036}.

\bibitem{Long2026}
C.~D. Long,
\emph{Counterexamples to the \(xz\)-Conjecture and the Mathieu
Conjecture for \(SU(2)\)},
arXiv:2607.19012 (2026).
\href{https://arxiv.org/abs/2607.19012}{arXiv:2607.19012}.

\bibitem{Dvorsky2026}
A.~Dvorsky,
\emph{Comment on ``The new counterexample to the Jacobian conjecture''},
Secret Blogging Seminar, 23 July 2026.
\url{https://sbseminar.wordpress.com/2026/07/20/the-new-counterexample-to-the-jacobian-conjecture/}.

\end{thebibliography}

\end{document}
